% 本文件由 LaTeX 排版 Agent 根据有效 translation.json 自动生成。
% author=Theodoret; work_id=2704; title=三段论式的论证

\chapter{三段论式的论证}

% structure: p0001 source=h1 target=omitted reason=page-title-already-rendered-by-parent

% structure: p0002 source=h2 target=section reason=relative-heading-rank; base=h2

\section{\OriginalTerm{天主圣言}{God the Word}是\OriginalTerm{不变}{Immutable}的}

1. 我们宣认了圣父、圣子和圣神的一个\OriginalTerm{本质}{substance}，并一致同意它是\OriginalTerm{不变}{immutable}的。那么，如果\OriginalTerm{天主圣三}{Trinity}的本质是一个且是不变的，那么\OriginalTerm{独生子}{only-begotten}，祂是圣三的一个\OriginalTerm{位格}{person}，也是不变的。而且，如果祂是不变的，那么祂就不是通过\OriginalTerm{变化}{mutation}而\OriginalTerm{降生成人}{made flesh}，而是在取了肉身后被称为成了血肉。\par

2. 如果天主圣言是通过\OriginalTerm{变化}{mutation}变成血肉，那么祂就不是不变的。因为没有一个有理智的人会把经历\OriginalTerm{变化}{alteration}的东西称为不变的。而如果祂是可变的，那么祂就与生祂者不是同一本质。因为一个\OriginalTerm{单纯本质}{uncompounded substance}的一部分可变而另一部分不变，这怎么可能呢？如果我们承认这一点，我们就会陷入\Person{阿里乌}{Arius}和\Person{欧诺米}{Eunomius}的\OriginalTerm{亵渎}{blasphemy}，他们断言圣子是\OriginalTerm{异质}{another substance}。\par

3. 如果主与圣父是\OriginalTerm{同体}{consubstantial}的，而圣子是通过变成血肉而成为血肉的，那么本质就同时是\OriginalTerm{可变}{mutable}和\OriginalTerm{不变}{immutable}的，这种\OriginalTerm{亵渎}{blasphemy}如果有人胆敢坚持，无疑会因对圣父的亵渎而加重，因为圣父分享同一本质，他肯定会称祂为可变的。\par

4. 在\OriginalTerm{圣经}{divine Scriptures}中记载，天主圣言\OriginalTerm{取了肉身}{took flesh}，也取了\OriginalTerm{灵魂}{soul}。最神圣的\Person{圣史}{Evangelist}说：圣言成了血肉。因此，我们必须在以下两者中择一：要么我们承认圣言\OriginalTerm{变化}{mutation}成血肉，从而拒绝所有新旧圣经，认为它们教导谎言；要么我们顺从圣经，承认肉身的\OriginalTerm{摄取}{assumption}，从思想中排除变化，以虔诚的心对待圣史的话。我们必须选择后者，因为我们宣认天主圣言的\OriginalTerm{本性}{nature}是不变的，并且有无数为摄取肉身作证的见证。\par

5. 居住在\OriginalTerm{帐幕}{tabernacle}中的不同于所居住的帐幕。圣史称肉身为帐幕，并说天主圣言\OriginalTerm{居住}{dwelt}在其中。他说：\BibleQuote{圣言成了血肉，住在我们中间。}\BibleRef{若 1:14}现在，如果祂通过变化成了血肉，祂就不是居住在肉身中。但我们受到的教导是祂居住在肉身中；因为同一圣史在另一处称祂的\OriginalTerm{身体}{body}为\OriginalTerm{圣殿}{temple}。因此，我们必须相信圣史对那些看似模糊之事的解释和诠释。\par

6. 如果圣史在写下\BibleQuote{圣言成了血肉}\BibleRef{若 1:14}之后，没有加上任何能消除\OriginalTerm{模糊}{ambiguous}的话，那么关于这段经文的\OriginalTerm{争论}{controversy}或许由于所用术语的隐晦而有一些合理的借口。但既然他紧接着说\BibleQuote{并住在我们中间}\BibleRef{若 1:14}，那么\OriginalTerm{争辩者}{combatants}的争论就是徒劳的。前一句由后一句得到解释。\par

7. 天主圣言的\OriginalTerm{不变性}{immutability}由最智慧的\Person{圣史}{Evangelist}明确宣告，因为他说\BibleQuote{圣言成了血肉，住在我们中间}\BibleRef{若 1:14}之后，立即加上\BibleQuote{我们看见了祂的光荣，正如父独生者的光荣，充满恩宠和真理。}\BibleRef{若 1:14}但是，如果按照愚人的说法，祂经历了\OriginalTerm{变化}{mutation}成了血肉，祂就不会保持祂原来的所是；然而，即使当祂被肉身包裹时，祂仍然放射出祂父尊贵的荣光，因此祂所拥有的\OriginalTerm{本性}{nature}是不变的，并且即使在身体内也闪耀着，传播那不可见本性的光辉。因为没有什么能遮蔽那\OriginalTerm{光明}{light}。\BibleQuote{光明在黑暗照耀，黑暗不能胜过它}\BibleRef{若 1:5}，正如极其神圣的若望所说。\par

8. 杰出的圣史渴望解释\OriginalTerm{独生子}{only-begotten}的\OriginalTerm{光荣}{glory}，但无法实现他的目的。因此他通过祂与圣父的\OriginalTerm{共融}{fellowship}来显示它。因为他说祂是属于那\OriginalTerm{本性}{nature}的；正如一个人看到\Person{若瑟}{Joseph}处于与其身份不符的奴役中，并且不知道他出身的高贵，便指出他的父亲是\Person{雅各伯}{Jacob}，先祖是\Person{亚巴郎}{Abraham}。同样，在这个意义上，圣史说当祂住在我们中间时，并没有减弱祂本性的光荣，\BibleQuote{因为我们看见了祂的光荣，正如父独生者的光荣。}\BibleRef{若 1:14}所以，即使当祂成了血肉时，祂是谁也是明显的，那么祂保持了原来的所是，并没有经历变成血肉的变化。\par

9. 我们宣认天主圣言不仅取了\OriginalTerm{身体}{body}，也取了\OriginalTerm{灵魂}{soul}。那么为什么神圣的圣史在此处省略了灵魂的提及，而只提肉身呢？这岂不是他展示了可见的本性，并借此指示了与之结合的本性？因为提到肉身当然也包括了灵魂的理解。因为当我们听到\Person{先知}{prophet}说\BibleQuote{凡有血肉的都要赞美祂的圣名}\BibleRef{咏 145:21}时，我们并不理解为先知是在劝勉没有灵魂的肉体，而是相信通过部分召唤，全体都被召唤来赞美。\par

10. 这些话\BibleQuote{圣言成了血肉}\BibleRef{若 1:1-2,14}显然不是指变化，而是指祂不可言喻的\OriginalTerm{慈爱}{loving-kindness}。因为杰出的圣史在说了\BibleQuote{在太初有圣言，圣言与天主同在，圣言就是天主}\BibleRef{若 1:1}，并宣告祂是可见与不可见之物的\OriginalTerm{创造者}{Creator}，称祂为\OriginalTerm{生命}{life}和\OriginalTerm{真光}{true light}，加上其他类似的表达，并以人类理性所能接受和语言所能表达的方式谈论了\OriginalTerm{神性}{Godhead}之后，他接着说\BibleQuote{圣言成了血肉}\BibleRef{若 1:14}，仿佛为无限的慈爱所震惊和惊奇。祂的存在是永恒的；祂是天主；祂创造了万有；祂是永恒生命和真光的源头；为了人类的\OriginalTerm{救恩}{salvation}，祂披上了肉身的\OriginalTerm{帐幕}{tabernacle}。人们以为祂只是所显现的那个样子。因此，圣史甚至没有提到\OriginalTerm{灵魂}{soul}，只提到可朽和\OriginalTerm{可朽的}{mortal}的肉身。关于\OriginalTerm{不朽的}{immortal}的灵魂，他什么也没有说，以显明慈爱的无限。\par

11. 神圣的\Person{宗徒}{Apostle}称主基督为\OriginalTerm{亚巴郎的后裔}{seed of Abraham}。但如果这是真的，而它确实是真，那么天主圣言就不是\OriginalTerm{变化}{changed}成了肉身，而是\OriginalTerm{摄取}{took on}了亚巴郎的\OriginalTerm{后裔}{seed}，按照宗徒本人的\OriginalTerm{教导}{teaching}。\par

12. 天主曾向达味发誓，要从他腰间的果实，按肉身，兴起基督，正如先知所说，又如同伟大的伯多禄所解释的。但如果圣言在变成肉身之后才被称为基督，那么我们在誓言中就找不到真理。然而我们被教导：天主不能撒谎，反而祂自己就是真理。因此，圣言并没有变成肉身，而是按照应许，取了达味后裔的初果。\par

% structure: p0015 source=h2 target=section reason=relative-heading-rank; base=h2

\section{证明联合不混淆}

1. 那些相信联合之后天主性与人性只有一性的人，藉此推论破坏了各自本性的特性；而本性的破坏意味着对其中任何一性的否认。因为联合本性的混淆妨碍我们认识肉身是肉身，天主是天主。但如果即使在联合之后，联合本性的区别仍然清晰，那就意味着没有混淆，联合是不混淆的。如果这一点被承认，那么主基督就不是一性，而是一位显示两性完好无损的圣子。\par

2. 我们也确认联合，并亲自承认联合发生在成孕之时；那么，如果通过联合，本性被混合和混淆，为何在出生之后，肉身没有显现任何新的特质，而是展现了人的特征，保持了婴孩的尺寸，被包裹在襁褓中，吮吸母亲的乳房？如果这一切不是仅仅发生在幻想和表象中，那么他们既不承认幻想也不承认表象；那么所见的确是一个身体。如果这一点被承认，那么本性并没有被联合所混淆，而是各自保持完好无损。\par

3. 这种拼凑而不协调的异端的倡导者，有时断言圣言成了肉身，有时又宣称肉身变成了天主性的本性。两种说法都是徒劳、空虚且充满谎言的，因为如果圣言如他们所主张的成了肉身，那么他们为什么只称祂为天主，并且拒绝也称祂为人，还严厉责备我们这些除了承认祂是天主之外也称祂为人的人？但如果肉身变成了天主性的本性，那么他们为什么用身体的预像来代替？因为当实体被毁灭时，预像是多余的。\par

4. 无形的本性不会在身体上受割损，但\Quote{在身体上}一词是因着心灵的精神割损而补充的；所以割损属于一个身体；而主基督在联合之后受了割损。如果这一点被承认，那么混淆的论证就被驳倒了。\par

5. 我们得知救主基督曾饥饿、口渴，我们相信这确实发生而非表象，但这样的状态不属于无形的本性，而属于身体。因此，主基督有一个身体，在复活之前按其本性感受影响。天主的宗徒为此作证，说：\BibleQuote{因为我们所有的，不是一位不能同情我们弱点的大司祭，而是一位在各方面受过试探的，与我们相似，只是没有罪过。}\BibleRef{希 4:15} 因为罪不是本性的，而是邪恶意志的。\par

6. 关于天主性，达味先知说：\BibleQuote{看，那保护以色列的，不打盹，也不睡眠。}\BibleRef{咏 121:4} 但福音的记载描述主基督在船上睡觉。如今不打盹与睡眠是两个相反的概念，因此如果如他们所说主基督只是天主，先知就与福音矛盾了。但并无矛盾，因为预言和福音都出自同一圣神。因此主基督有一个身体，与所有其他身体相似，受睡眠需要的支配。所以混淆的论证被证明是虚构的。\par

7. 关于天主性，依撒意亚先知说：\BibleQuote{他不饥饿，也不疲倦，}\BibleRef{依 40:28} 等等。但福音说：\BibleQuote{耶稣因旅途疲惫，就坐在井旁。}\BibleRef{若 4:6} 而\Quote{不疲倦}与\Quote{疲倦}相反。因此预言与福音的记载矛盾。但它们并不矛盾，因为两者都出自同一天主。不疲倦属于不受限制的本性，它充满万有。但从一处移动到另一处属于受限制的本性；当那移动者被迫旅行时，它承受旅客的疲劳。因此行走并疲倦的是一个身体，因为联合并没有混淆各本性。\par

8. 主基督对被囚禁在狱中的圣保禄说：\BibleQuote{不要害怕，保禄，}\BibleRef{宗 18:9} 等等。但同一位基督，祂驱走了保禄的恐惧，自己却如此害怕，以致如圣路加所证，全身汗如血滴，洒满周围的土地，并藉天使的援助得到力量。这些陈述彼此对立，因为害怕怎能与驱走恐惧相反？然而它们并不对立。因为同一位基督按其本性是天主也是人；作为天主，祂坚固需要安慰的人；作为人，祂通过天使接受安慰。尽管天主性和圣神作为傅油同在，但那时身体和灵魂并未得到与它们联合的天主性或圣神的支持，而是将这项服务委托给一位天使，为要显明灵魂和身体的软弱，并藉着软弱使人看到软弱的各本性。这些事显然因天主性的许可而发生，为要使将来生活在后世的人中，那些相信灵魂和身体被摄取的人能因这些证明而得证，而反对者则被明确的证据驳倒。如果联合是通过成孕实现的，并且如他们所说使两性成为一性，那么本性的属性如何能保持完好无损，灵魂痛苦，身体因极度恐惧而汗如血滴？但如果一种是身体的本性，另一种是灵魂的本性，那么联合并没有造成肉身和天主性的一性，而是显示了一位圣子，在祂内同时显出人性和天主性。\par

9. 倘使他们说，复活后身体变成了天主性，则可恰当地如此回答：即使在复活后，身体仍被看见有限定地有手足和一切肢体部分；它是可触摸、可见的；它有伤口和疤痕，如同复活前一样。因此必须坚持两者之一：要么这些部分归于神性本体——如果身体在变为神性本性时仍具有这些部分；要么反之，必须承认身体仍在其本性的界限内。然而，神性本体是单纯而不复合的，而身体是复合且分为许多部分的；所以它并未变成天主性的本性，即使在复活后虽为不朽、不腐、充满神圣荣光，它仍是有其自身限定的身体。\par

10. 主在复活后，向不信的宗徒们显示了他的手、脚和钉痕；然后为了教导他们所见并非幻象，他补充说：\BibleQuote{鬼神是没有骨肉的，你们看我，却是有的。}因此，身体并未变成神体，而是有骨有肉、有手有脚。所以，即使在复活后，身体仍是身体。\par

11. 神性本体是不可见的，但三倍有福的斯德望说他看见了主，所以即使在复活后，主的身体仍是身体，并被得胜的斯德望看见了，因为神性本体是不能被看见的。\par

12. 如果所有人类都将看见人子乘着天上的云彩来临，依照主自己的话，而他对梅瑟说过：\BibleQuote{没有人看见了我还能活着。}两者皆真，那么他将带着他升天时的身体而来。因为那身体是可见的，天使也对宗徒们说过：\BibleQuote{这位离开你们、被接到天上去的耶稣，你们看见他怎样升了天，也要怎样降来。}如果这是真的——诚然是真的——那么肉体和神性并非同一本性，而是无混淆的结合。\par

% structure: p0028 source=h2 target=section reason=relative-heading-rank; base=h2

\section{证明救主的神性不受苦难}

1. 我们既由圣经和集于尼西亚的圣教父们所教导，宣认圣子与天主圣父同一性体。父的不受苦性也由本性所教导并为圣经所宣告。那么我们将进一步宣认圣子也是不受苦难的，因为这一界定因同一性体而被强化。因此，每当我们听见圣经宣告主基督的十字架和死亡时，我们将苦难归于肉身，因神性本体按其本性是不受苦的，绝不能受苦。\par

2. 主基督说：\BibleQuote{凡父所有的，都是我的。}其中一样便是不受苦性。因此，如果作为天主，他是不受苦的，那么他便作为人而受苦。因为神性本体不受苦。\par

3. 主说：\BibleQuote{我所要赐给的食粮，就是我的肉，是为世界的生命而赐给的。}又说：\BibleQuote{我是善牧，我认识我的羊，我的羊也认识我……我为羊舍掉我的性命。}所以，身体和灵魂都由善牧为那有灵魂和身体的羊而给出。\par

4. 人的本性由身体和灵魂合成。但人犯罪了，需要无任何瑕疵的祭献。因此造物主取了身体和灵魂，使它们远离罪污，为人的身体给出他的身体，为人的灵魂给出他的灵魂。如果这是真的——诚然是真的，因为这是真理之言——那么那些将苦难归于神性本体的人是狂妄而亵渎的。\par

5. 有福的保禄称基督为\BibleQuote{死者中的首生者}；我想首生者与他被称为首生者的那些人有相同的本性。因此，作为人，他是死者中的首生者，因为他首先摧毁了死亡的痛苦，并给众人带来了另一生命的甜美希望。他如何复活，也就如何受苦。因此，作为人他受苦，但作为可敬畏的天主，他是不受苦的。\par

6. 神圣的宗徒称我们的救主基督为\BibleQuote{亡者的初果}，但初果与它所出自的整体有关系。因此，他不是作为天主而被称为初果，因为天主性与人性之间有何关系？前者是不朽的本性，后者是可朽的。这便是那些睡眠之人的本性，基督被称为他们的初果。死亡和复活属于这一本性，而在他复活中，我们有了一般复活的证据。\par

7. 当主基督愿说服不信的宗徒他已消灭死亡并复活时，他向他们显示了他身体的部分——他的肋旁、他的手、他的脚，以及其中保存的苦难痕迹。那么，这身体复活了，并且我相信，这是向不信者显示的。复活的便是被埋葬的，被埋葬的便是死了的，死了的当然是被钉在十字架上的。所以与身体结合的神性本体是不受苦的。\par

8. 那些描述主的肉身为生命赐予者的人，以自己的话使生命本身成为可朽的。他们应当明白，肉身是通过与之结合的生命而成为生命赐予者的。但如果按他们的说法，生命是可朽的，那么肉体本身按本性是可朽的，并通过生命成为生命赐予者，它如何能保持生命赐予呢？\par

9. 天主圣言按本性是不朽的，而肉体按本性是可朽的，但在苦难之后，通过与圣言的结合，肉体本身成为不朽的。那么，说这样不朽的赐予者分享了死亡，岂不荒谬？\par

10. 那些坚持天主圣言在肉体中受苦的人，应当被问及他们所说之话的含义；如果他们大胆回答说，当身体被钉子刺穿时，神性本体感受到了痛苦，那么让他们知道，神性本体并未担任灵魂的角色。天主圣言随同身体也取了灵魂。如果他们拒绝这一论证为亵渎，并断言肉体按本性受苦，而天主圣言将这苦难视为他自己肉身的苦难，那么让他们不要提出令人困惑和晦涩的短语，而要清楚地说出那刺耳短语的含义。所有愿意追随圣经的人，在这解释中都会支持他们。\par

11. 宗徒之长伯多禄在其公函中说，基督在肉身内受了苦。但听见基督受苦的人，并不以为无体的天主圣言受苦，而是以为降生成人的那一位受苦。基督的名字指明两种性体；但\Quote{肉身}一词与苦难相连，并非指明两者受苦，而是指明两者之一受苦。因为听见基督在肉身内受苦的人，会认为祂既是天主，便是不能受苦的，于是将苦难只归于肉身。正如我们听见他说，天主曾向达味起誓，要从他腰间的果实按肉身兴起基督，我们并不说天主圣言源自达味，而是说天主圣言所取的肉身与达味有血缘关系；同样，听见基督在肉身内受苦的人，也必须承认苦难属于肉身，并宣认天主性的不能受苦。\par

12. 当主基督在十字架上说：\Quote{父啊！我把我的灵魂交在你手中}，这灵魂被亚略派和欧诺弥派说成是独生子的天主性，因为他们认为祂所取的肉身是没有灵魂的；但真理的宣讲者却说，这灵魂正是如此称呼的，其根据来自以下经文。明智的福音作者立即接着说：\Quote{说了这话，便断了气}。路加如此说，真福马尔谷也同样接着说：\Quote{便断了气}。神圣的玛窦写道：\Quote{交付了灵魂}；神圣的若望写道：\Quote{交付了灵魂}。众人皆按人的习惯说法，因为我们惯常用这些说法来描述死亡的人；这些说法无一含有天主性的意义，而是全指灵魂。若有人接受亚略派对这段经文的解释，即便如此，它仍显示天主性的不死性。因为基督将它交托给父，并没有将它交予死亡。那么，那些否认灵魂被摄取、坚持天主圣言是受造物、并声称祂在肉身内代替灵魂的人，既然否认祂被交付于死亡，他们又如何能获得宽恕呢？而那些承认圣三同性体，并让灵魂保留其不死性的人，竟厚颜无耻地敢说与圣父同性体的天主圣言尝了死亡，他们又如何能得赦免呢？\par

13. 如果基督既是天主又是人，正如神圣圣经所教导、卓越教父们一贯宣讲的，那么祂作为人受苦，作为天主则不能受苦。\par

14. 如果承认肉身的摄取，并宣告它在复活前是可受苦的，又宣讲天主性是不能受苦的，那么为何要离开那可受苦的性体，而将苦难归于那不能受苦的性体呢？\par

15. 如果我们的主和救主将字据钉在十字架上，如神圣宗徒所说，那么祂钉住的便是肉身，因为每个人都在自己身上像文字一样刻下自己罪恶的痕迹；因此，祂为罪人舍了那完全无罪的身体。\par

16. 当我们说身体、肉身或人性受苦时，我们并没有分离天主性；因为正如天主性曾与那饥饿、口渴、疲倦、甚至睡眠、并经历苦难的人性结合，它本身不受这一切影响，却容许人性按自身方式受影响；同样，当祂被钉十字架时，天主性也与之结合，并容许苦难的完成，以便通过苦难摧毁死亡。它并未从苦难中感受到痛苦，却将苦难当作自己的苦难，视之为自己的圣殿和与自己结合之肉身的苦难；正因为这肉身，信友们被称为基督的肢体，而祂自己也被称为信者们的头。\par

\SourceRecord{Translated by Blomfield Jackson.；Nicene and Post-Nicene Fathers, Second Series；Vol. 3.；Edited by Philip Schaff and Henry Wace.；Buffalo, NY: Christian Literature Publishing Co.,；1892.；<http://www.newadvent.org/fathers/2704.htm>.}
